% Source printed pp. 4-12, entire Section 1.1.
\section{模理论简史}\label{src:1.1}
\begin{epigraph}
Riemann 的精神将像激励我们一样, 激励未来的世代.
\epigraphauthor{Lars Ahlfors, \cite[第 $53$ 页]{Ahl53}}
\end{epigraph}
模理论的历史发展, 使我们得以初步窥见模理论中的许多主题.

\subsection{Bernhard Riemann 与 \texorpdfstring{$M_g$}{Mg} 的起源}\label{src:1.1.1-subsection}
\begin{epigraph}
\textit{Die $3p-3$ übrigen Verzweigungswerthe in jenen Systemen gleichverzweigter $\mu$ werthiger Functionen können daher beliebige Werthe annehmen; und es hängt also eine Klasse von Systemen gleichverzweigter $\overline{2p+1}$ fach zusammenhängender Functionen und die zu ihr gehörende Klasse algebraischer Gleichungen von $3p-3$ stetig veränderlichen Grössen ab, welche die Moduln dieser Klasse genannt werden sollen.}

译文: 这些分歧方式相同的 $\mu$ 值函数系统中, 余下的 $3p-3$ 个分支点因而可以取任意给定的值; 于是, 一类具有相同分歧方式, 连通度为 $2p+1$ 的函数系统, 以及相应的一类代数方程, 依赖于 $3p-3$ 个连续变量, 我们将它们称为这一类的模.
\epigraphauthor{Bernhard Riemann, \cite[第 $33$ 页]{Rie57}}
\end{epigraph}
这是一篇非凡论文中的一句非凡的话: Riemann 既引入了``模''的概念, 又计算出 $M_g$ 的``模的数目''是 $3g-3$. Riemann 的思路大致如下: 不考虑抽象的光滑曲线, 而把曲线看成 $\mathbb P^1$ 上的分歧覆盖, 并考虑模空间
\begin{equation}\label{src:1.1.1}
\Hur_{d,g}=\left\{[C\to\mathbb P^1]\;\middle|\;
\begin{aligned}
&\text{$C$ 是亏格为 $g$ 的光滑曲线}\\
&\text{$C\to\mathbb P^1$ 是次数为 $d$ 的单分歧覆盖}
\end{aligned}\right\}.
\end{equation}
Hurwitz 后来正式研究了这些模空间 \cite{Hur91}; 它们如今称为 \newterm{Hurwitz 空间}{Hurwitz space}, 在研究 $M_g$, 尤其是在不可约性的证明中, 起着关键作用(见第~\bookref{6.7}~节).

\newterm{单分歧覆盖}{simply branched covering}是光滑曲线之间的有限映射, 其分歧指数至多为 $2$, 且每根纤维至多含有一个分歧点; 另见定义~\bookref{6.7.2}. 由 Riemann--Hurwitz 公式~\bookeqref{6.7.4}, 每个单分歧覆盖 $C\to\mathbb P^1$ 都在 $\mathbb P^1$ 上 $2g+2d-2$ 个互异的点处分支. 这给出交换图
\begin{equation}\label{src:1.1.2}
\begin{tikzcd}[column sep=2.8em,row sep=2.2em]
& \Hur_d(C)\arrow[r,hook]\arrow[dl] & \Hur_{d,g}\arrow[dl]\arrow[dr] &\\
{[C]}\arrow[r,hook] & M_g & & \Sym^{2d+2g-2}\mathbb P^1
\end{tikzcd}
\end{equation}
其中
\begin{itemize}
\item 映射 $\Hur_{d,g}\to\Sym^{2d+2g-2}\mathbb P^1$ 把覆盖 $[C\to\mathbb P^1]$ 送到这 $2g+2d-2$ 个分支点; 这里 $\Sym^N\mathbb P^1=(\mathbb P^1)^N/S_N$ 是分类 $N$ 个无序点的空间; 
\item 映射 $\Hur_{d,g}\to M_g$ 定义为 $[C\to\mathbb P^1]\mapsto[C]$; 
\item $\Hur_d(C)$ 是 $[C]\in M_g$ 在映射 $\Hur_{d,g}\to M_g$ 下的原像, 即 $\Hur_d(C)$ 分类曲线 $C$ 固定时的单分歧覆盖 $C\to\mathbb P^1$.
\end{itemize}
如果 $d$ 充分大, 那么对 $\mathbb P^1$ 上任意一般的 $2d+2g-2$ 个点的集合, 都存在一条亏格为 $g$ 的曲线 $C$ 和一个单分歧覆盖 $C\to\mathbb P^1$, 它恰好在这些点处分支; 而且这样的映射至多有有限多个. 换言之, $\Hur_{d,g}\to\Sym^{2d+2g-2}\mathbb P^1$ 的像稠密, 且纤维有限; 精确细节见第~\bookref{6.7.2-subsection}~小节. 因此, 
\begin{equation}\label{src:1.1.3}
\begin{aligned}
\dim M_g&=\dim\Hur_{d,g}-\dim\Hur_d(C)\\
&=2d+2g-2-\dim\Hur_d(C).
\end{aligned}
\end{equation}
为计算 $\Hur_d(C)$ 的维数, 注意一个单分歧覆盖 $C\to\mathbb P^1$ 由一个次数为 $d$ 的线丛 $L$(即 $\Pic^d(C)$ 中的一个元素)和两个无基点的截面给出, 使得诱导的到 $\mathbb P^1$ 的映射是单分歧的. 由于一般选取的两个截面定义单分歧覆盖(引理~\bookref{6.7.16}), 我们得到
\begin{equation}\label{src:1.1.4}
\dim\Hur_d(C)=\dim\Pic^d(C)+2h^0(C,L)-1,
\end{equation}
这里减去 $1$, 是因为同时缩放两个截面会定义相同的到 $\mathbb P^1$ 的映射. Riemann--Roch 定理~\bookeqref{6.1.5}~告诉我们 $h^0(C,L)=d+1-g$. 另一方面, $\dim\Pic^d=\dim\Pic_0=g$; 这可以利用指数正合列看出: $0\to\mathbb Z\to\mathcal O_C\xrightarrow{\exp}\mathcal O_C^*\to0$ 给出长正合列
\begin{equation}\label{src:1.1.5}
\underbrace{H^1(C,\mathbb Z)}_{\mathbb Z^{2g}}
\longrightarrow\underbrace{H^1(C,\mathcal O_C)}_{\mathbb C^g}
\longrightarrow\underbrace{H^1(C,\mathcal O_C^*)}_{\Pic(C)}
\xrightarrow{\deg}\underbrace{H^2(C,\mathbb Z)}_{\mathbb Z},
\end{equation}
并给出等同 $\Pic_0(C)\cong\mathbb C^g/\mathbb Z^{2g}$. 因此, $\dim\Pic_0(C)=g$, 利用~\eqref{src:1.1.4}, 算得 $\dim\Hur_d(C)=g+2(d+1-g)-1=2d-g+1$. 将其代入~\eqref{src:1.1.3}, 得到
\[
\dim M_g=(2d+2g-2)-(2d-g+1)=3g-3.
\]
事实上, Riemann 还给出了几个计算 $M_g$ 维数的其他启发式论证. 关于 $M_g$ 的模的数目的进一步讨论, 见 \cite[第 $255$--$257$ 页]{GH78} 或 \cite[第 $211$--$215$ 页]{Mir95}; 关于 Riemann 这些计算的历史背景, 见 \cite{AJP16}.
\begin{quote}
\textbf{Riemann 的模问题: }$M_g$ 是否作为复解析空间存在?
\end{quote}
虽然如今可以用现代方法使 Riemann 的论证完全严格(本书后面也会这样做), 但 Riemann 的方法存在基础性问题; 今天我们会说, Riemann 计算的是一个``局部形变空间''的维数. 尤其是, 当时并不知道 $M_g$ 是否存在, 也不清楚 $M_g$ 应该是哪一种空间. 尽管如此, Riemann 已经直觉地把握了它的几何; 事实上, 在同一篇论文 \cite{Rie57} 中, 他引入了 ``Mannigfaltigkeit''(即``流形'')一词来描述它的几何. 直到很久以后, 在 Teichmüller, Chern 和 Weil 的工作之后, 流形才于 $1940$ 年代得到正式定义.

\subsection{低亏格曲线的模}\label{src:1.1.2-subsection}
\begin{epigraph}
\textit{我认为, 这个簇很可能是有理的, 或者至少能与线性空间中点组的一个对合建立对应...}
\epigraphauthor{Francesco Severi 的猜想: 对所有 $g$, $M_g$ 都是单有理的, \cite[第 $881$ 页]{Sev15}.}
\end{epigraph}
\paragraph*{亏格 $0$.} 当 $n\geq3$ 时, 带有 $n$ 个有序互异点的光滑亏格 $0$ 曲线的模空间 $M_{0,n}$ 可以描述为
\[
M_{0,n}=\left(\mathbb P^1\setminus\{0,1,\infty\}\right)^{n-3}\setminus(\text{所有对角线}).
\]
事实上, 给定 $\mathbb P^1$ 上 $n$ 个有序互异点 $p_1,\ldots,p_n$, 存在唯一的自同构 $g\in\Aut(\mathbb P^1)\cong\PGL_2$, 将 $(p_1,p_2,p_3)$ 送到 $(0,1,\infty)$. 当 $n=4$ 时, 得到双射 $M_{0,4}=\mathbb P^1\setminus\{0,1,\infty\}$, 由 $\mathbb P^1$ 上四个点的经典交比给出. 交比最早由亚历山大的 Pappus 于公元 $300$ 年发现 \cite{Ale86}; 另见例~\bookref{8.2.8}.

\paragraph*{亏格 $1$.} Dedekind 和 Klein 已经知道椭圆曲线的模. 每条椭圆曲线 $(E,p)$, 即带有一个标记点 $p\in E$ 的光滑亏格 $1$ 曲线 $E$, 都可以描述为 Legendre 形式的三次平面曲线
\[
E_\lambda=V(y^2z-x(x-z)(x-\lambda z))\subseteq\mathbb P^2,
\]
其中 $\lambda\neq0,1$, 且 $p=[0:1:0]\in E_\lambda$. 但 $\lambda$ 的选择并不唯一: $\lambda$, $1/\lambda$, $1-\lambda$, $1/(1-\lambda)$, $\lambda/(\lambda-1)$ 和 $(\lambda-1)/\lambda$ 给出同构的椭圆曲线. 换言之, 由 $\lambda\mapsto[E_\lambda]$ 给出的映射 $\mathbb A^1\setminus\{0,1\}\to M_{1,1}$ 是一个 $6$ 对 $1$ 的满射. 另一方面, $j$ 不变量
\[
j=2^8\frac{(\lambda^2-\lambda+1)^3}{\lambda^2(\lambda-1)^2}
\]
唯一确定曲线的同构类, 因而给出双射 $M_{1,1}\cong\mathbb A^1$. 现代论述见 \cite[\S$4$]{Har77}.

\paragraph*{亏格 $2$.} 每条光滑亏格 $2$ 曲线 $C$ 都是超椭圆的, 可以写成 $\mathbb P^1$ 上的二重覆盖 $y^2=(x-a_1)\cdots(x-a_6)$. 这是因为微分层 $\Omega_C$ 是次数为 $2$, 具有 $2$ 个整体截面的无基点线丛; 由 Riemann--Hurwitz 公式~\bookeqref{6.7.4}, 诱导的映射 $C\to\mathbb P^1$ 在 $6$ 个点处分歧. 得到描述
\[
M_2=\left(\Gamma(\mathbb P^1,\mathcal O(6))\setminus\Delta\right)/\GL_2,
\]
其中 $\Delta\subseteq\Gamma(\mathbb P^1,\mathcal O(6))$ 表示具有重根的二元六次式构成的轨迹. 在 $\mathbb P^1$ 上作射影坐标变换, 可以使曲线在 $0,1,\infty$ 及另外 $3$ 个点 $a_4,a_5,a_6\in\mathbb P^1\setminus\{0,1,\infty\}$ 处分歧. 这样便得到满射 $(\mathbb P^1\setminus\{0,1,\infty\})^3\setminus\Delta\to M_2$.

二元六次式的不变量理论(见 \cite{Cle70})给出更精确的描述: 不变多项式环, 即二元六次式系数中在 $\mathbb P^1$ 的自同构下不变的多项式所成的环, 由不变量 $J_2,J_4,J_6,J_8,J_{10},J_{15}$ 生成, 下标表示次数, 并满足唯一的关系 $J_{15}^2=G(J_2,J_4,J_6,J_{10})$, 其中 $G$ 为某个多项式. 不变量 $J_{10}$ 是二元六次式的判别式, 而 $J_{15}$ 不影响概形结构. 由此, $M_2$ 是加权射影空间中的开子集
\[
M_2=\Proj\mathbb C[J_2,J_4,J_6,J_{10}]\setminus\{J_{10}=0\},
\]
这说明 $M_2$ 是一个通过下列函数嵌入 $\mathbb A^8$ 的仿射簇: 
\[
\frac{J_2^5}{J_{10}},\quad
\frac{J_2^3J_4}{J_{10}},\quad
\frac{J_2^2J_6}{J_{10}},\quad
\frac{J_2J_4^2}{J_{10}},\quad
\frac{J_2J_6^3}{J_{10}^2},\quad
\frac{J_4^5}{J_{10}^2},\quad
\frac{J_4J_6}{J_{10}},\quad
\frac{J_6^5}{J_{10}^3}.
\]
可以把这个坐标环等同于 $\mathbb Z/5$ 在 $\mathbb A^3$ 上作用的不变量环, 其中生成元 $\zeta\in\mathbb Z/5$ 的作用为 $\zeta\cdot(x,y,z)=(\zeta x,\zeta^2y,\zeta^3z)$; 上述函数分别对应不变量 $x^5$, $x^3y$, $x^2z$, $xy^2$, $xz^3$, $y^5$, $yz$, $z^5$. 这给出相当优美的整体描述
\[
M_2=\mathbb A^3/(\mathbb Z/5).
\]
Bolza 曾在经典工作中研究这一点 \cite{Bol87}, 较近的研究则见 Igusa~\cite{Igu60}.

\paragraph*{亏格 $3$.} 非超椭圆的光滑亏格 $3$ 曲线通过典范嵌入成为 $\mathbb P^2$ 中的一条四次曲线. 令 $\Delta\subseteq\Gamma(\mathbb P^2,\mathcal O(4))$ 为奇异四次曲线的轨迹, 则 $M_3$ 中非超椭圆曲线构成的开轨迹可描述为商 $(\Gamma(\mathbb P^2,\mathcal O(4))\setminus\Delta)/\GL_3$; 这是第一次, 我们的描述只能描述一般的曲线. 另一方面, 超椭圆的亏格 $3$ 曲线是 $\mathbb P^1$ 上在 $8$ 个点处分歧的二重覆盖, 于是得到集合意义下的分解
\[
M_3=\underbrace{\left(\Gamma(\mathbb P^2,\mathcal O(4))\setminus\Delta\right)/\GL_3}_{\dim=6}
\coprod\underbrace{\left(\Gamma(\mathbb P^1,\mathcal O(8))\setminus\Delta\right)/\GL_2}_{\dim=5},
\]
这提示超椭圆曲线的轨迹是 $M_3$ 中的一个除子.

\paragraph*{亏格 $4$.} 非超椭圆的光滑亏格 $4$ 曲线 $C$ 通过典范嵌入到 $\mathbb P^3$, 并可实现为一个二次曲面 $Q$ 与一个三次曲面 $S$ 的交 $C=Q\cap S$. 这给出有理映射 $\Gamma(\mathbb P^3,\mathcal O(2))\times\Gamma(\mathbb P^3,\mathcal O(3))\dashrightarrow M_4$, 其像是非超椭圆曲线的轨迹; 同上, 超椭圆轨迹可由 $(\Gamma(\mathbb P^1,\mathcal O(10))\setminus\Delta)/\GL_2$ 参数化. 另一种方式是把一般的非超椭圆光滑亏格 $4$ 曲线实现为恰有两个结点的五次平面曲线的正规化, 或实现为 $\mathbb P^1$ 上在 $12$ 个点处分支的三重覆盖.

\paragraph*{亏格 $5$--$10$.} 经典上, 低亏格曲线或者通过映射 $C\to\mathbb P^2$ 的像, 被描述为具有指定奇点的平面曲线, 或者被描述为分歧覆盖 $C\to\mathbb P^1$. 对一般的亏格 $g$ 曲线 $C$, 使 $C$ 可实现为奇异平面曲线的正规化的最小次数是 $d=\lfloor(2g+8)/3\rfloor$. 如果该平面曲线最坏只有结点奇点, 那么结点的数目为 $\delta:=(d-1)(d-2)/2-g$. 同时, 映射 $C\to\mathbb P^1$ 的最小次数为 $\lfloor(g+3)/2\rfloor$. 另见 \cite[第 $21$ 页]{Mum75a}.
\begin{table}[htbp]
\centering\small
\begin{tabular}{@{}cccccc@{}}
\toprule
$g$ & \shortstack{$d=\operatorname{im}(C\to\mathbb P^2)$\\的最小次数} & $\delta=$ 结点数 & $\dfrac{(d+1)(d+2)}2-3\delta$ & \shortstack{$C\to\mathbb P^1$\\的最小次数} & 分支点数\\\midrule
$0$ & $1$ & $0$ & $3$ & $1$ & $0$\\
$1$ & $3$ & $0$ & $10$ & $2$ & $4$\\
$2$ & $4$ & $1$ & $12$ & $2$ & $6$\\
$3$ & $4$ & $0$ & $15$ & $3$ & $10$\\
$4$ & $5$ & $2$ & $15$ & $3$ & $12$\\
$5$ & $6$ & $5$ & $13$ & $4$ & $16$\\
$6$ & $6$ & $4$ & $16$ & $4$ & $18$\\
$7$ & $7$ & $8$ & $12$ & $5$ & $22$\\
$8$ & $8$ & $13$ & $6$ & $5$ & $24$\\
$9$ & $8$ & $12$ & $9$ & $6$ & $28$\\
$10$ & $9$ & $18$ & $1$ & $6$ & $30$\\
$11$ & $10$ & $25$ & $-9$ & $7$ & $34$\\\bottomrule
\end{tabular}
\caption{一般的低亏格曲线.}\label{src:1.1.6}
\end{table}
在 \cite{Sev15,Sev21} 中, Severi 使用这些描述证明了当 $g\leq10$ 时 $M_g$ 是单有理的. 如同他那个时代的其他数学家, Severi 并未精确表述 $M_g$ 作为模空间究竟意味着什么.

当 $g\geq11$ 时, 问题出在哪里?随着亏格增大, 描述一般的亏格 $g$ 曲线也变得更加困难. 为说明 $g\geq11$ 时的困难, 尝试把亏格 $g$ 曲线描述为次数 $d=\lfloor(2g+8)/3\rfloor$, 在指定点 $p_1,\ldots,p_\delta\in\mathbb P^2$ 处具有 $\delta=(d-1)(d-2)/2-g$ 个结点的平面曲线. 若 $f\in\Gamma(\mathbb P^2,\mathcal O(d))$ 定义的平面曲线在每个 $p_i$ 处有结点, 则方程 $f_x(p_i)=f_y(p_i)=f_z(p_i)=0$ 在 $\Gamma(\mathbb P^2,\mathcal O(d))$ 上施加 $3\delta$ 个线性方程. 要使这样的结点平面曲线存在, 需要
\[
\dim\Gamma(\mathbb P^2,\mathcal O(d))-3\delta
=\frac{(d+1)(d+2)}2-3\left(\frac{(d-1)(d-2)}2-g\right)>0.
\]
如表~\ref{src:1.1.6}~所示, $g=11$ 是这一条件首次不成立的情形!

\paragraph*{Severi 猜想.} 虽然 Severi 关于所有 $g$ 的 $M_g$ 都是单有理的猜想最终被证明是错误的, 但它仍激励了数学家数十年: ``是否还有更多 $M_g$, $g\geq11$, 是单有理的, 是一个非常有趣的问题, 但它看来也很困难, 尤其当 $g$ 很大时'' \cite[第 $37$ 页]{Mum75a}. 在 $1980$ 年代, Eisenbud, Harris 和 Mumford 推翻了这一猜想, 证明在大亏格时, 某种意义下恰恰相反: 当 $g\geq24$ 时, $M_g$ 是一般型的 \cite{HM82,EH87}.

\paragraph*{Petri 对典范曲线的描述.}
\begin{epigraph}
〔Petri 的方法〕不可避免地有些繁琐, 但即使只是为了能夸口, 我也觉得能说``我曾见过\underline{每一条}曲线''是一件好事.
\epigraphauthor{David Mumford, \cite[第 $17$ 页]{Mum75a}}
\end{epigraph}
$19$ 世纪和 $20$ 世纪初的大多数数学家把曲线描述为奇异平面曲线, 或描述为 $\mathbb P^1$ 的覆盖, 而 Petri 对典范嵌入曲线的显式描述 \cite{Pet23} 是一个例外, 更接近现代方法. M. Noether 证明了非超椭圆光滑曲线 $C$ 的典范嵌入 $C\hookrightarrow\mathbb P^{g-1}$ 是射影正规的 \cite{Noe80}, 即映射 $\varphi:\Sym^*\Gamma(C,\Omega_C)\to\bigoplus_{d\geq0}\Gamma(C,\Omega_C^{\otimes d})$ 是满射. 在此结果以及 Enriques~\cite{Enr19} 和 Babbage~\cite{Bab39} 的工作基础上, Petri 证明: 齐次理想 $I=\ker\varphi$ 由二次式生成, 除非 $C$ 是五次平面曲线($g=6$)或三角曲线(即 $\mathbb P^1$ 的三重覆盖); 在这两种例外情形下, $I$ 由次数为 $2$ 和 $3$ 的式子生成. Petri 的分析具有显著的构造性, 给出了在 $\mathbb P^{g-1}$ 中截出 $C$ 的显式方程, 以及这些方程之间的显式合冲. Petri 的工作至今仍启发着模理论和合冲理论中的研究. 本书不再进一步介绍这一观点, 但推荐 \cite{SD73}, \cite[第 $17$--$21$ 页]{Mum75a}, \cite{AS78,Gre82,Gre84} 和 \cite[\S III.$3$]{ACGH85}.

\subsection{解析方法与 Teichmüller 空间}\label{src:1.1.3-subsection}
\begin{epigraph}
一件作品以其优美令你喜爱, 并不必然意味着创造它的人就值得你的尊敬.
\epigraphauthor{Plutarch, \textit{名人传}}
\end{epigraph}
在 $19$ 世纪末和 $20$ 世纪初, 人们将 Riemann 曲面描述为上半平面关于 $\mathrm{PSL}_2(\mathbb R)$ 的离散子群的商; 这些子群以 Fuchs~\cite{Fuc66} 命名, 称为 \newterm{Fuchs 群}{Fuchsian group}. Fricke 和 Klein 在共 $1300$ 页的著作 \cite{FK92,FK12} 中, 使用自守函数理论对 Fuchs 群进行分类. 他们构造了今天称为 Teichmüller 空间的空间, 证明它可缩, 甚至给出了复结构. Torelli 证明了 Riemann 曲面可以从其 Jacobian 重建 \cite{Tor13}, 而 Siegel 将维数为 $g$ 的 Abel 簇的模空间 $A_g$ 构造为一个解析空间 \cite{Sie35}.

Oswald Teichmüller 首次精确表述了 Riemann 的模问题, 把 $M_g$ 构造为复解析空间, 并将 $3g-3$ 解释为其复维数 \cite{Tei40,Tei44}. Teichmüller 构造了 Teichmüller 空间 $T_g$, 参数化亏格为 $g$ 的拓扑曲面 $\Sigma_g$ 上按同胚等同的复结构. 空间 $T_g$ 同胚于 $\mathbb C^{3g-3}$ 中的一个球, 并带有映射类群 $\Gamma_g$ 的作用; $\Gamma_g$ 是 $\Sigma_g$ 的微分同胚群关于与恒等映射同痕的微分同胚所成子群的商. 这一作用是正常不连续的, $M_g$ 可实现为商 $T_g/\Gamma_g$. 虽然 Teichmüller 理论曾被遗忘近 $20$ 年, 后来 Ahlfors, Bers, Weil 等人极大地扩展了它; 见 \cite{Wei57,Wei58,AB60,Ber60,Ahl61}. 现代的介绍性论述见 \cite{Ber72,Hub06,FM12}.

Teichmüller 还引入了 Riemann 曲面族的概念, 并证明 Teichmüller 空间满足泛性质. Grothendieck 在 Cartan 讨论班的十次讲座 \cite{Gro61} 中, 用范畴和函子的语言发展了解析模空间的一般理论, 在这一框架中重新表述了 Teichmüller 理论, 并证明 $T_g$ 表示参数化 Riemann 曲面族的函子. 这为 Grothendieck 后来关于代数模的工作奠定了基础: ``可以期待, 有一天我们能够从 Teichmüller 空间理论中完全消除分析, 因为它应当是纯粹几何的'' \cite[讲座 I]{Gro61}.

\subsection{代数模理论的起源}\label{src:1.1.4-subsection}
\paragraph*{Boole, Cayley, Gordan 和 Hilbert.} $19$ 世纪的不变量理论家关心的是, 在射影自同构意义下分类 $n$ 个变量的 $d$ 次齐次多项式; 换言之, 研究模空间 $\Gamma(\mathbb P^{n-1},\mathcal O(d))/\PGL_n$. 他们试图通过给出多项式 $f=\sum_Ia_Ix^I$ 的系数 $a_I$ 中的显式不变多项式, 来描述这个模空间. 不变量理论起源于 George Boole~\cite{Boo41} 和 Arthur Cayley~\cite{Cay45} 的工作, 并由 Paul Gordan, David Hilbert 及许多其他人进一步发展. Gordan 给出了二元型($n=2$)的不变量环的显式生成元 \cite{Gor68}, Hilbert 后来证明了对任意环, 不变量环都是有限生成的 \cite{Hil90,Hil93}. 我们将在推论~\bookref{7.5.13}(3)~和注记~\bookref{7.5.2}~中证明并讨论 Hilbert 的定理.

Cayley 构造了 $\mathbb P^3$ 中曲线的模空间 \cite{Cay60,Cay62}, 如今称为 Chow 簇. 他的想法是给次数为 $d$ 的曲线 $C\subseteq\mathbb P^3$ 对应上与 $C$ 有非空交的直线 $L\subseteq\mathbb P^3$ 的集合. 这是 $\mathrm{Gr}(1,\mathbb P^3)$ 中的一个 $d$ 次超曲面, 而 Chow 簇是在 $\Gamma(\mathrm{Gr}(1,\mathbb P^3),\mathcal O(d))$ 中所有这样得到的超曲面构成的集合的闭包. Chow 和 van der Waerden 后来将此构造推广到 $\mathbb P^n$ 中任意维数的子簇 \cite{CW37}; 见第~\bookref{2.4.5-subsection}~小节. Chow 簇将在第~\bookref{6.8}~节中构造 $\overline M_g$ 时发挥作用.

\paragraph*{Weil.} André Weil 在研究有限域上曲线的 Riemann 假设 \cite{Wei48} 时, 需要构造曲线的 Jacobian, 以参数化次数为 $0$ 的线丛. 当时, 人们只考虑嵌入在仿射空间或射影空间中的簇; 在其奠基性工作 \cite{Wei62} 中, Weil 将这个范畴扩大到抽象簇. 这足以构造 Jacobian, 并证明曲线的 Riemann 假设---事实上, 这是他的第二个证明. 后来, Weil 和 Chow 独立证明了 Jacobian 是射影的.

\paragraph*{Baily.} Walter Baily 将主极化 Abel 簇的模空间 $A_g$ 构造为拟射影簇 \cite{Bai60a,Bai60b}, 证明了 Satake 的拓扑紧化 \cite{Sat56} 是代数的 \cite{Bai58}, 并与 Borel 一起引入了今天称为 \newterm{Baily--Borel 紧化}{Baily--Borel compactification} 的构造 \cite{BB66}. 利用把曲线送到其 Jacobian 的周期映射 $M_g\to A_g$, 以及断言这一映射单射的 Torelli 定理, Baily 得出 $M_g$ 具有拟射影簇结构. 但他没有证明这给出的是``自然''的簇结构, 也没有证明它具有任何唯一性, 即没有证明 $M_g$ 是一个粗模空间.
\begin{quote}
第三, 要把 $\mathcal E$ 称为亏格 $n$ 的 Riemann 曲面的模簇, 就应当能够断言, 它是唯一的, 并且在某种意义下, 在亏格 $n$ 曲线的代数系统的正规参数簇中具有普遍性. 也就是说, 给定任意亏格 $n$ 曲线的正规代数系统(这里指其参数簇为正规簇), 该系统中非奇异成员的参数簇应当有一个到 $\mathcal E$ 的自然映射.\hfill---Baily, \cite[第 $59$--$60$ 页]{Bai60b}
\end{quote}
Mumford 在 \cite[第 $98$ 页]{Mum75a} 中将 $M_g$ 的拟射影性归功于 Baily, Gieseker 在 \cite{Mum04} 的评注中也持相同看法.

\paragraph*{Grothendieck.} Alexander Grothendieck 在 \cite{Gro61} 中将解析模理论形式化之后, 在 ``FGA 系列'' \cite{FGAI,FGAIII,FGAIV,FGAV,FGAVI} 中把函子方法应用于代数几何: 他引入 Hilbert, Quot 和 Picard 函子, 证明它们可由射影概形表示, 发展下降理论, 并引入预叠和叠的概念. 当然, Grothendieck 后来还通过发展概形理论, 重新建立了代数几何的全部基础. 他对代数几何乃至更广泛的数学的深刻影响, 塑造了模理论的未来.

尽管 Grothendieck 没有发表关于 $M_g$ 的论文, 但他仍十分关心 $M_g$ 作为拟射影簇的存在, 以及它在任意特征下的连通性; 这可从他在 $1960$ 年代初与 Mumford 等人的通信中看出 \cite[\S II]{Mum10}. Grothendieck 意识到, 自同构的存在阻碍了参数化光滑曲线族的函子的可表性. 他通过同时参数化曲线上的 $n$ 级结构, 即 $H_1(C,\mathbb Z/n\mathbb Z)$ 的一组辛基, 来刚性化模问题. 虽然他能够证明带有 $n\geq3$ 级结构的光滑曲线函子可由概形表示, 但证明其拟射影性却颇为困难. 他的想法是对作用在 $n$ 级结构选择上的有限群取商, 把 $M_g$ 构造为刚性化模空间的商. 缺乏拟射影性阻碍了这一方法, 因为一个非拟射影簇关于有限群的商未必作为簇存在.

\paragraph*{Mumford.} 受到 Riemann 模空间, 以及 $20$ 世纪初 Chow 簇, Picard 簇和 Abel 簇的模的构造的启发, David Mumford 对模理论的基础作出了巨大贡献, 并首次系统研究其几何. 通过结合 Grothendieck 的概形理论形式体系与 $19$ 世纪的不变量理论, Mumford 发展了如今称为几何不变量理论(Geometric Invariant Theory, 简称 GIT)的代数几何商理论, 我们将在第~\bookref{8}~章介绍它. Mumford 应用 GIT 构造了 $M_g$ 和 $A_g$. 其理论最初在 \cite{Mum61} 中作了概述, 并在奠基性著作 \cite{GIT} 中完整展开. 事实上, Mumford 给出了 $\Spec\mathbb Z$ 上粗模概形 $M_g$ 的两种构造 \cite[定理 $5.11$ 和 $7.13$]{GIT}, 并证明对于每个素数 $p$, $M_g$ 在 $\Spec\mathbb Z[1/p]$ 上是拟射影的.\footnote{有趣的是, Mumford 的两种构造实际上都没有使用 GIT, 至少没有使用通常被视为``标准 GIT 工具''的那套方法, 即通过 Hilbert--Mumford 判别法验证 GIT 稳定性. Mumford 的一种构造依赖于主极化 Abel 簇的模空间 $A_g$ 的存在, 另一种则依赖于使用协变式的特殊方法.} $\overline M_g$ 在 $\Spec\mathbb Z$ 上的射影性, 后来由其他方法证明 \cite{Knu83b,Mum77,Gie82}, 这些方法更直接适用于其他模空间.

Mumford 还构造了一个参数化固定光滑曲线上稳定向量丛的拟射影簇 \cite{Mum63}; Seshadri 随后证明, 半稳定向量丛的模空间给出了射影紧化 \cite{Ses67}.

在奠基性的合作论文 \cite{DM69} 中, Deligne 和 Mumford 引入了稳定曲线及 $\mathcal M_g$ 的紧化 $\overline{\mathcal M}_g$, 证明了稳定约化定理~\bookeqref{6.5.1}, 并首次引入代数叠的概念, 如今称为 Deligne--Mumford 叠. 最后, 他们给出了任意特征下 $M_g$ 连通性的两个证明\footnote{虽然两个证明都根本性地使用了紧化 $\overline{\mathcal M}_g$ 和稳定约化定理, 但代数叠的使用并不是本质的.}; 见第~\bookref{6.7}~节. 遗憾的是, Deligne 和 Mumford 没有证明他们关于叠的任何断言: ``本节结果的证明将在别处给出'' \cite[第 $76$ 页]{DM69}. 严格基础的缺失, 使代数叠在此后数十年中获得了令人望而生畏的名声. 直到 $2000$ 年代, 代数叠才进入代数几何的主流. 在 Deligne 和 Mumford 的论文发表后的 $50$ 多年中, 其大多数断言如今都已在文献中得到完整证明. 现在已有介绍代数叠的优秀教材 \cite{LMB00,Ols16}, 而 Johan de Jong 的 Stacks Project~\cite{SP} 则提供了无可置疑的坚实基础.

\paragraph*{Artin.} 代数空间与叠的理论由 Michael Artin 发展. 类似于 Weil 将仿射簇和射影簇扩大到抽象簇, 将概形范畴扩大到代数空间, 使我们能够构造有限群作用的商, 更一般地, 构造任意平展等价关系的商.\footnote{Matsusaka 也通过考虑某些等价关系的商, 建立了 Q-簇理论 \cite{Mat64}, 但它不如代数空间理论完善.} Artin 的学生 Knutson 首次系统写出了代数空间理论 \cite{Knu71}.

在两篇论文 \cite{Art69a,Art69b} 中, Artin 证明了模理论中的两个关键结果: Artin 逼近定理~\bookeqref{B.5.18}~和 Artin 代数化定理~\bookeqref{C.6.8}. 在开创性的论文 \cite{Art74b} 中, Artin 引入了比 Deligne 和 Mumford 更广泛的代数叠概念, 包含 $\cBun_{r,d}(C)$ 这样的, 自同构群可能无限的叠. 他还给出了代数叠的局部形变论刻画, 称为 Artin 的代数性公理(\bookref{C.7.1}~和~\bookref{C.7.4}), 可用来验证给定的模叠是代数的; 见第~\bookref{C.7.2-subsection}~小节.

进一步的历史背景, 推荐阅读 \cite{Mum75a,Oor81,Kle05,JP13,AJP16,Kol21}.
